\documentclass[10pt,a4paper]{article}
\usepackage[T1]{fontenc}
\usepackage{lmodern}
\usepackage{microtype}
\usepackage[margin=25mm]{geometry}
\usepackage{amsmath,amssymb,amsthm}
\usepackage[colorlinks=true,linkcolor=blue,citecolor=blue,urlcolor=blue]{hyperref}
\hypersetup{pdftitle={Activation-error direction and matrix-output error}}

\newcommand{\R}{\mathbb{R}}
\newcommand{\norm}[1]{\left\lVert #1\right\rVert}
\newtheorem{proposition}{Proposition}

\title{Activation-error direction and matrix-output error}
\author{}
\date{15 November 2026}

\begin{document}
\maketitle
\noindent\textbf{Status:} mathematical observation; established linear algebra.
No novelty or empirical performance claim is made.

\section{Statement and scope}
Two activation errors of the same magnitude can affect the receiving matrix
multiplication very differently: their contributions can reinforce or cancel.
The relevant quantity depends on both the error vector and the receiving
weights, rather than on activation error magnitude alone.

Let a fixed affine operation be
\begin{equation}
  y=Wx+b,\qquad W\in\R^{m\times n},\quad x\in\R^n,\quad b\in\R^m.
  \label{aed:operation}
\end{equation}
An approximation, such as quantization, replaces $x$ by $\widehat{x}=x+e$.
Compare outputs with the same $W$ and $b$ in exact real arithmetic. The newly
introduced local output error is exactly
\begin{equation}
  d=\widehat{y}-y=W(\widehat{x}-x)=We.
  \label{aed:local-error}
\end{equation}
Here all vector norms are Euclidean. This isolates activation approximation:
weight changes, finite-precision accumulation, and error arriving from earlier
operations are not included in $d$.

\section{The receiving matrix defines an error geometry}
\begin{proposition}
Let $w_i$ denote column $i$ of $W$, and let $G=W^TW$. Then
\begin{align}
  \norm{d}_2^2
    &=e^TGe \label{aed:gram}\\
    &=\sum_{i=1}^n\norm{w_i}_2^2 e_i^2
      +2\sum_{i<j}(w_i^Tw_j)e_i e_j.
      \label{aed:cross-terms}
\end{align}
\end{proposition}
\begin{proof}
Substitute $d=We$ into $d^Td$ and expand the symmetric quadratic form.
\end{proof}

The diagonal terms measure individual contributions; the cross terms describe
their interaction. A negative cross term reduces the combined squared error,
while a positive one increases it. Even a weight-aware diagonal score
$\sum_i\norm{w_i}_2^2 e_i^2$ discards these interactions and is not, in
general, an upper bound.

Since $G$ is positive semidefinite,
\begin{equation}
  \lambda_{\min}(G)\norm{e}_2^2
  \leq\norm{We}_2^2
  \leq\lambda_{\max}(G)\norm{e}_2^2.
  \label{aed:spectral-bounds}
\end{equation}
Thus direction matters whenever $G$ has unequal eigenvalues. In particular,
an error in the null space of $W$ has zero local output effect, however large
its norm; an error of the same norm in a sensitive direction need not.
If $G=cI$, output error depends only on the input error norm instead.

\section{A quantization example with equal activation MSE}
Choose $W=\begin{bmatrix}1&1\end{bmatrix}$ and $b=0$. Use unit-step
round-to-nearest quantization $Q(t)=\operatorname{round}(t)$, without clipping.
Consider the two inputs
\begin{equation}
  x_+=\begin{bmatrix}1/4\\1/4\end{bmatrix},\qquad
  x_-=\begin{bmatrix}1/4\\-1/4\end{bmatrix}.
  \label{aed:example-inputs}
\end{equation}
Both quantize to zero. Their activation errors satisfy
\begin{equation}
  e_+=\begin{bmatrix}-1/4\\-1/4\end{bmatrix},\qquad
  e_-=\begin{bmatrix}-1/4\\1/4\end{bmatrix},\qquad
  \norm{e_+}_2^2=\norm{e_-}_2^2=\frac18.
  \label{aed:example-residuals}
\end{equation}
Their elementwise mean squared error is therefore $1/16$ in both cases.
Nevertheless,
\begin{equation}
  We_+=-\frac12,\qquad We_-=0,
  \qquad \norm{We_+}_2^2=\frac14,\quad\norm{We_-}_2^2=0.
  \label{aed:example-outputs}
\end{equation}
For $e_+$ the two contributions reinforce; for $e_-$ they cancel exactly.
The diagonal score is $1/8$ for both, so it underestimates the first output
error and overestimates the second. This is an exact counterexample to using
activation MSE alone to determine the receiving operation's output error.

\section{The same issue across finite matrix blocks}
Partition the reduction dimension into disjoint groups:
\begin{equation}
  W=\begin{bmatrix}W_1&\cdots&W_B\end{bmatrix},\qquad
  e=\begin{bmatrix}e_1^T&\cdots&e_B^T\end{bmatrix}^T,\qquad
  d_g=W_g e_g.
  \label{aed:blocks}
\end{equation}
Then
\begin{equation}
  d=\sum_{g=1}^B d_g,\qquad
  \norm{d}_2^2=\sum_g\norm{d_g}_2^2
                +2\sum_{g<h}d_g^Td_h.
  \label{aed:block-interactions}
\end{equation}
Adding squared block errors omits the interaction terms and is not a general
upper bound. A conservative bound that requires no independence assumption is
\begin{equation}
  \norm{d}_2\leq\sum_g\norm{d_g}_2.
  \label{aed:triangle-bound}
\end{equation}
The algebra applies to finite reduction blocks, but does not specify their
physical layout, scale handling, or accumulator arithmetic. A causal decision
can use observed contributions; it cannot rely on cancellation from unseen
future blocks.

\section{Interpretation and limits}
The example establishes a possibility, not how frequently useful cancellation
occurs in trained networks. Quantization errors are functions of current values;
zero mean, independence, and favorable directions must not be assumed without
evidence. Local output error also does not determine final task quality after
nonlinear operations or an iterative inference trajectory.

Related methods already exploit aspects of this problem: ResQ uses a local
norm-based surrogate for dynamic residual precision~\cite{resq}; AQuant adjusts
activation rounding to improve output error~\cite{aquant}; FGMP weights live
activation quantization differences with prepared sensitivity~\cite{fgmp}.
The identities above are reusable mathematical background, not a claim to
introduce those objectives or a new precision-selection algorithm.

\begin{thebibliography}{3}
\bibitem{resq}
D. Abati et al.
\emph{ResQ: Residual Quantization for Video Perception}.
ICCV, 2023, Sec.~4.2.
\url{https://arxiv.org/abs/2308.09511v1}.

\bibitem{aquant}
Z. Li et al.
\emph{Efficient Adaptive Activation Rounding for Post-Training Quantization}.
Author preprint, version 3, 2023, Secs.~3--4.
\url{https://arxiv.org/abs/2208.11945v3}.

\bibitem{fgmp}
C. Hooper et al.
\emph{FGMP: Fine-Grained Mixed-Precision Weight and Activation Quantization
for Hardware-Accelerated LLM Inference}.
Author preprint, version 1, 2025, Secs.~3--4.
\url{https://arxiv.org/abs/2504.14152v1}.
\end{thebibliography}
\end{document}
